\section*{Hoofdstuk \ref{ch:b3:microbiomes} --- Microbiomen en symbiosen}

\begin{solution}{exo:b3:microbiomes:1}
\omterm{def:b3:microbiomes:symbiosis}{Mutualisme}: beide winnen --- vlinderbloemige en rhizobium, koraal en
alg. \omterm{def:b3:microbiomes:symbiosis}{Commensalisme}: de een wint, de ander blijft onaangedaan --- de
meeste darmbacteriën die leven van wat de gastheer niet kan verteren.
Parasitisme: de een wint ten koste van de ander --- \emph{Wolbachia}
die mannelijke embryo's doodt. Hetzelfde paar kan verschuiven: een
darmcommensaal die de wand oversteekt wordt een ziekteverwekker, een
rhizobium dat ophoudt met fixeren wordt een parasiet van het knolletje,
en een alg onder hittestress schaadt het koraal dat haar huisvestte.
\end{solution}

\begin{solution}{exo:b3:microbiomes:2}
$H = -(0.5\ln 0.5 + 0.3\ln 0.3 + 0.1\ln 0.1 + 2\times 0.05\ln 0.05) =
0.347 + 0.361 + 0.230 + 0.300 = 1.24$; $e^{H} = 3.4$ effectieve
soorten. $D = 0.25 + 0.09 + 0.01 + 0.0025 + 0.0025 = 0.355$; $1/D =
2.8$.
\end{solution}

\begin{solution}{exo:b3:microbiomes:3}
Het 16S-onderzoek meldt wie er is, tot ongeveer het geslacht, en hun
relatieve aantallen \omterm{def:b3:genomics:ngs}{reads}; het zegt niets over de genen die zij dragen,
mist schimmels en virussen (andere \omterm{def:b3:genetic-engineering:pcr}{primers}), en is vertekend door het
aantal genkopieën en door de \omterm{def:b3:genetic-engineering:pcr}{primers}. De hagelschotmetagenomica meldt
genen, functies en, bij voldoende diepte, hele genomen en stammen; zij
kost meer per monster, wordt bij weefselmonsters overspoeld door
gastheer-DNA, en kan nog altijd niet zeggen welke genen tot expressie
komen of welke cellen leven.
\end{solution}

\begin{solution}{exo:b3:microbiomes:4}
Vezels vergisten tot kortketenige vetzuren die de dikke darm en de
gastheer voeden; vitamine~K en B-vitamines maken; galzuren en
geneesmiddelen omzetten; \omterm{def:b3:microbiomes:gut}{kolonisatieweerstand} tegen ziekteverwekkers;
het afweersysteem scholen. Ontwrichting: colitis door
\emph{Clostridioides difficile} na \omterm{def:b3:bacteriology:antibiotics}{antibiotica} (ook in verband
gebracht: chronische darmontsteking, obesitas).
\end{solution}

\begin{solution}{exo:b3:microbiomes:5}
$\binom{20}{5} = \num{15504}$. Soort met $10$: $1 - \binom{10}{5}/
\num{15504} = 1 - 252/\num{15504} = 0.984$; met $6$: $1 - 2002/\num{15504}
= 0.871$; met $3$: $1 - 6188/\num{15504} = 0.601$; met $1$: $1 -
\num{11628}/\num{15504} = 0.25$. $E[S_{5}] = 2.7$ soorten. Dekking: één
eenling, $1 - 1/20 = 0.95$.
\end{solution}

\begin{solution}{exo:b3:microbiomes:6}
$400\times 10^{11} = 4\times 10^{13}$ bacteriën, ongeveer \qty{40}{g}.
Genen: $1000$ soorten $\times$ \num{4000}, minus de overlap --- in de
orde van $3\times 10^{6}$ verschillende genen, meer dan honderd keer de
menselijke \num{20000}.
\end{solution}

\begin{solution}{exo:b3:microbiomes:7}
Het \omterm{def:b3:microbiomes:nodule}{nitrogenase} wordt door zuurstof vernietigd, maar de bacteroïden
moeten ademen om zestien ATP per N$_{2}$ te maken. De schors van het
knolletje vormt een diffusiebarrière die de intrede van zuurstof
beperkt, en het \omterm{def:b3:microbiomes:nodule}{leghemoglobine}, in millimolaire concentratie, bindt de
zuurstof die binnenkomt zodat de vrije concentratie bij
\qty{10}{nM} blijft terwijl de stroom naar de bacteroïden groot is.
Roze is oxyleghemoglobine; een groen knolletje heeft zijn
\omterm{def:b3:microbiomes:nodule}{leghemoglobine} afgebroken --- het is verouderd en fixeert niet meer.
\end{solution}

\begin{solution}{exo:b3:microbiomes:8}
Beide falen: geen knolletjes en geen \omterm{def:b3:microbiomes:mycorrhiza}{arbusculaire mycorrhiza}, omdat de
calciumpiek het gedeelde signaal van de \omterm{def:b3:microbiomes:mycorrhiza}{gemeenschappelijke symbioseroute}
is. Die route gaat dus aan de knolvorming vooraf --- zij is
vierhonderd miljoen jaar geleden voor de schimmelsymbiose gebouwd en
later door de vlinderbloemigen voor de \omterm{def:b3:microbiomes:nodule}{rhizobia} ingelijfd.
\end{solution}

\begin{solution}{exo:b3:microbiomes:9}
Kiers stelde sommige knolletjes van soja bloot aan een atmosfeer van
argon en zuurstof, waarin de \omterm{def:b3:microbiomes:nodule}{rhizobia} geen stikstof konden binden,
terwijl de overige knolletjes van de plant normaal fixeerden; de plant
kneep de zuurstoftoevoer naar de niet-fixerende knolletjes af en hun
\omterm{def:b3:microbiomes:nodule}{rhizobia} plantten zich ongeveer half zo sterk voort. Zonder dat
onderscheid zouden \omterm{def:b3:microbiomes:nodule}{rhizobia} die de kostbare fixatie oversloegen zich
sneller voortplanten dan de fixeerders, zich door de populatie
verspreiden, en zouden de planten eindigen met bacteriën die niets
gaven: het \omterm{def:b3:microbiomes:symbiosis}{mutualisme} zou instorten.
\end{solution}

\begin{solution}{exo:b3:microbiomes:10}
$\binom{N-N_{i}}{n}/\binom{N}{n} = \prod_{j=0}^{n-1}(N - N_{i} - j)/(N -
j)$, met elke factor onder $1$; van $n$ naar $n + 1$ gaan
vermenigvuldigt met nog een factor onder $1$, zodat de kans om te
missen daalt en de kans op aanwezigheid met $n$ stijgt; bij $n = N$ is
de teller $\binom{N -
N_{i}}{N}$ nul en is elke soort aanwezig, wat $S$ geeft. Voor $N_{i},
n \ll N$ is elke factor ongeveer $1 - N_{i}/N$ en het product ongeveer
$(1 -
N_{i}/N)^{n} \approx e^{-nN_{i}/N} \approx (1 - n/N)^{N_{i}}$ --- de
benadering van trekken met teruglegging.
\end{solution}

\begin{solution}{exo:b3:microbiomes:11}
Een mannetje is een doodlopende weg voor een symbiont die alleen in
eieren reist, dus verhoogt elke eigenschap die mannetjes in vrouwtjes
verandert, ze doodt ten gunste van hun zusters, of niet-besmette
vrouwtjes benadeelt, de overdracht van de symbiont. Cytoplasmatische
onverenigbaarheid: het sperma van besmette mannetjes wordt zo gewijzigd
dat de zygote sterft tenzij het ei dezelfde \emph{Wolbachia} draagt, die
haar redt. Niet-besmette vrouwtjes verliezen dus het nageslacht dat zij
met besmette mannetjes verwekken, terwijl besmette vrouwtjes er geen
verliezen; het voordeel groeit met de besmettingsfrequentie, zodat de
besmetting boven een drempel tot fixatie doorzet, zelfs als besmette
vrouwtjes iets minder eieren leggen --- de kost weegt niet op tegen het
deel van de nesten van niet-besmette vrouwtjes dat wordt vernietigd.
\end{solution}

\begin{solution}{exo:b3:microbiomes:12}
De endosymbiont gaat bij elke gastheergeneratie door een flessenhals
van enkele cellen en recombineert nooit met buitenstaanders: de drift
legt licht schadelijke mutaties vast, deleties inbegrepen, en er is
geen bron van vervangende genen (de ratel van Muller). Genen waarvan de
gastheer de producten levert, of die een vrij leven dienen ---
omhulsel, beweeglijkheid, regulatie, het meeste herstel --- worden niet
gemist en gaan het eerst; de translatiemachinerie en de genen die maken
wat de gastheer nodig heeft (de essentiële aminozuren bij
\emph{Buchnera}) blijven het langst. Vrijlevende verwanten hebben
enorme populaties, recombinatie en een instroom van genen, en de
selectie houdt hun genomen in stand. Omkering vergt nieuwe genen van
buiten --- in afzondering onmogelijk --- zodat gastheren in plaats
daarvan een uitgesleten symbiont door een verse vervangen, zoals
verscheidene insectenlijnen hebben gedaan.
\end{solution}

\begin{solution}{pb:b3:microbiomes:1}
\textbf{1.} $400\times 10^{11} = 4\times 10^{13}$ bacteriën, \qty{40}{g}.
\textbf{2.} Een genoom met \num{4000} genen is ongeveer \qty{4}{Mb},
\qty{4.4}{fg}: $4\times 10^{13}\times 4.4\times 10^{-15} = \qty{0.18}{g}$
bacterieel DNA, tegen $3\times 10^{13}\times 6.4\times 10^{-12} = \qty{0.19}{g}$
menselijk DNA --- ongeveer evenveel.
\textbf{3.} $1000\times 0.7\times 4000 + 0.3\times 4000 \approx 2.8\times
10^{6}$ verschillende genen, zo'n $140$ keer dat van het menselijk
genoom.
\textbf{4.} $30\times 10 = \qty{300}{mmol}$; $\times 0.9 = \qty{270}{kJ}$;
\qty{3}{\%} van de voeding.
\textbf{5.} Nee --- \qty{3}{\%} verklaart geen \qty{30}{\%}. Kiemvrije
muizen nemen ook minder doeltreffend op, hebben andere darmhormonen,
vetopslag en warmteproductie, en een andere eetlust: het \omterm{def:b3:microbiomes:symbiosis}{microbioom}
werkt op de fysiologie van de gastheer en niet alleen als
brandstofleverancier.
\textbf{6.} Deling in evenwicht met lozing: een stabiele populatie
waarvan elke cel dagelijks wordt vervangen. Na een doding van
\qty{99.9}{\%} blijven er $4\times 10^{10}$ over en herstellen tien
verdubbelingen de aantallen binnen ongeveer tien dagen --- maar de
overlevenden en de snelste groeiers overheersen, de samenstelling is
veranderd (een \omterm{def:b3:microbiomes:gut}{dysbiose}), en in dat gat kan een indringer als
\emph{C.~difficile} zich vestigen.
\textbf{7.} De kleine soorten delen $0.4$ over $177$: $p = 0.00226$ elk. $H
= -(0.3\ln 0.3 + 0.2\ln 0.2 + 0.1\ln 0.1 + 0.4\ln 0.00226) = 0.361 +
0.322 + 0.230 + 2.437 = 3.35$; $e^{H} \approx 28$.
\textbf{8.} $D = 0.09 + 0.04 + 0.01 + 177\times (0.00226)^{2} = 0.141$;
$1/D = 7.1$. De \omterm{def:b3:microbiomes:diversity}{simpsonindex} kwadrateert de abundanties, zodat de drie
overheersende soorten haar bepalen en de zeldzame nauwelijks meetellen;
Shannon geeft zeldzame soorten meer gewicht.
\textbf{9.} $1 - 60/2000 = 0.97$: ongeveer $30$ \omterm{def:b3:genomics:ngs}{reads} per duizend
behoren tot soorten die nog niet zijn gezien.
\textbf{10.} De rijkdom stijgt met de monstergrootte terwijl er zeldzame
soorten blijven verschijnen; \num{10000} \omterm{def:b3:genomics:ngs}{reads} bereiken soorten die
\num{2000} missen. Rarefieer het grote monster met de stelling (of door
te deelsteekproeven) tot \num{2000} \omterm{def:b3:genomics:ngs}{reads} --- het zou ongeveer $180$
moeten geven als het dezelfde gemeenschap is --- en vergelijk op
gelijke diepte.
\textbf{11.} $N_{i} = 1$: $1 - \binom{1999}{1000}/\binom{2000}{1000} =
1 - (2000 - 1000)/2000 = 0.5$. $N_{i} = 5$: ongeveer $1 - 0.5^{5} = 0.97$.
\textbf{12.} $e^{2.0} = 7.4$ effectieve soorten; $D \ge 0.5^{2} = 0.25$,
ongeveer $0.27$: een gemeenschap met één overheersende soort en enkele
tientallen kleine --- de ingestorte darm met weinig diversiteit waarin
\emph{C.~difficile} vat krijgt.
\textbf{13.} $150\times 8 = \qty{1200}{kg}$ suiker per hectare,
\qty{10}{\%} van de \qty{12}{t} die is gefotosynthetiseerd.
\textbf{14.} $150\,000/28 = 5360$ mol N$_{2}$; $16\times 5360 =
\num{85700}$ mol ATP.
\textbf{15.} $\num{85700}/30 = 2860$ mol glucose, \qty{514}{kg} ---
ongeveer \qty{43}{\%} van de betaalde suiker. De rest bouwt en
onderhoudt de knolletjes en de bacteroïden, levert de acht elektronen
per N$_{2}$ als reductiemiddel, en bouwt de ammoniak in.
\textbf{16.} $1200\times 16 = \qty{19200}{MJ}$ voor \qty{150}{kg} N:
\qty{128}{MJ/kg}, bijna vier keer de \qty{35}{MJ/kg} van de fabriek ---
maar betaald in zonlicht, op het veld, zonder brandstof.
\textbf{17.} Fixeerders $0.9\times 1 = 0.9$; profiteurs $0.1\times 0.5 =
0.05$; de profiteurs zijn $0.05/0.95 = 5.3\,\%$ --- in één seizoen
gehalveerd.
\textbf{18.} Profiteurs $0.1\times 1.2^{10} = 0.62$ tegen fixeerders
$0.9$: \qty{41}{\%} na tien seizoenen, op weg naar overheersing.
Sancties keren de selectie tegen de profiteur; zonder hen brokkelt het
\omterm{def:b3:microbiomes:symbiosis}{mutualisme} af.
\textbf{19.} $10^{6}\times \qty{20}{pg} = \qty{20}{\micro g}$ koolstof
per cm$^{2}$ per dag; \qty{18}{\micro g} naar het koraal.
\textbf{20.} Ademhaling \qty{15}{\micro g}: een overschot van
\qty{3}{\micro g} per cm$^{2}$ per dag voor groei, de organische matrix
van de kalkafzetting, slijm en gameten.
\textbf{21.} Met \qty{5}{\%} van de algen \qty{0.9}{\micro g} per dag
tegen \qty{15}{\micro g} verademd: een tekort van \qty{14}{\micro g} per
dag; \qty{300}{\micro g} voorraad houdt ongeveer drie weken.
\textbf{22.} Vier weken op $+2$ is erger: de schade aan de
fotosystemen van de algen stijgt steil en niet lineair met de
temperatuur. De maat werkt niettemin als waarschuwing omdat de
verbleking de opgetelde stress samenvat en de twee verlopen minder
verschillen dan de onzekerheid in de drempel; het is een empirische
index, geijkt op waargenomen verbleking.
\textbf{23.} Een koraal onder stress kan verschuiven naar
warmtebestendiger algensoorten die al in kleine aantallen aanwezig
zijn, of ze na de verbleking uit het water opnemen; de bestendige algen
geven minder koolstof, zodat het koraal trager groeit, en de gewonnen
bestendigheid is een of twee graden --- minder dan de voorspelde
opwarming --- terwijl het tempo van de opwarming de aanpassing van de
koralen zelf voorbijstreeft.
\textbf{24.} $10^{10}$ cm$^{2}$ $\times$ \qty{20}{\micro g} $=
\qty{200}{kg}$ koolstof per dag, ongeveer \qty{73}{t} per jaar.
\textbf{25.} Ongeveer \qty{3}{\%} van de voeding uit de vergisting van
het \omterm{def:b3:microbiomes:symbiosis}{microbioom}; \qty{97}{\%} dekking van het sequencingmonster;
\qty{10}{\%} van de fotosynthese betaald voor de stikstof van het
bonengewas.
\end{solution}
